\chapter{Conclusion et perspectives}
\label{chap:conclusion}

Le stage a permis de faire évoluer une base électronique existante vers un dossier de fabrication, d'établir une première architecture firmware et de construire une chaîne progressive de validation. Après un banc exploratoire STM32/IMU, le simulateur de mission formalise le levé coopératif de cinq AUV, l'absence de radio sous l'eau, les rendez-vous de surface, la supervision, le protocole compact et les scénarios de sécurité. Sa conteneurisation et son déploiement Web en font également un démonstrateur autonome. Le HIL a ensuite exécuté cinq échanges bidirectionnels authentifiés entre un ESP32 et un Raspberry Pi réels, avec un modèle logiciel du lien nRF24.

Ce qui subsiste au-delà du stage ne se limite pas à un dossier de fabrication. Le simulateur, sa suite de tests et l'implémentation portable du protocole constituent des briques réutilisables par l'équipe : elles décrivent le comportement attendu d'une flotte Caiman indépendamment d'une carte particulière, et fournissent une référence contre laquelle le firmware définitif pourra être confronté. Le banc HIL, en exécutant ce même code C sur deux processeurs distincts, établit le pont entre cette référence et le monde physique ; il restera valable lorsque le transport Wi-Fi sera remplacé par le lien nRF24 réel.

La principale limite est l'absence de réception du PCB avant la rédaction. Aucun résultat de bring-up, de mesure d'alimentation, de capteur ou de radio sur la carte Caiman ne peut donc être revendiqué. La chaîne de vérification s'arrête à la CAO et à la revue de fabrication : elle établit que le dossier est cohérent et fabricable, non que la carte fonctionne.

Les prochaines étapes prioritaires en découlent directement : la réception et l'inspection visuelle des cartes, le contrôle de continuité des rails avant toute mise sous tension, la vérification des marges des inductances, la montée en tension avec limitation de courant, la mesure des rails, la programmation du STM32 puis le test progressif de chaque périphérique. L'intégration du protocole HIL au firmware de la carte permettra enfin de relier les résultats numériques au système physique.
